\itemhead{1.2 思考 · 未编号题组（教材p.16）}{XJ-B1-C01-TH-P016-Y0927}
\label{XJ-B1-C01-TH-P016-Y0927}
{\footnotesize 湘教版教材p.16 / PDF第23页 / 1.2 思考 / 未编号题组}
\field{原题}{对于 1.2.1 小节例 2 后面列出的 6 个命题，哪些是真命题？哪些是假命题？
\par\textit{所引用的教材 p.15 六个命题如下：}
\subq{1}若两个三角形全等，则它们相似；
\subq{2}若两个三角形相似，则它们全等；
\subq{3}若实数 $a\ne0$，则 $a^2>0$；
\subq{4}若四边形 $ABCD$ 为菱形，则 $AC\perp BD$；
\subq{5}若 $a<0$，则方程 $ax^2+4x-3=0$ 没有正的实根；
\subq{6}若 $a=b$，则 $a^2=b^2$。}
\field{参考答案}{依次为真、假、真、真、假、真。（2）反例为大小不同的等边三角形；（5）取 $a=-1$，方程有正根1、3。}
{\small \field{关联知识}{L01 命题、真假、逆命题与否定（教材帮p.30 / PDF39）}
\field{方法与易错点}{明确论域与断言；真命题给理由，假命题给满足题设而使结论失败的反例。 易错：把逆命题当成否定；只改变结论中的个别词而改变原逻辑范围。}
\field{答案依据}{答案册未提供本栏目的独立对应答案；依据原题、教材概念补写参考回应。开放性总结不设唯一数值答案。}}
